\documentclass[journal]{IEEEtran}
\usepackage{amsmath,amsfonts}
\usepackage{algorithmic}
\usepackage{algorithm}
\usepackage{array}
\usepackage[caption=false,font=normalsize,labelfont=sf,textfont=sf]{subfig}
\usepackage{textcomp}
\usepackage{stfloats}
\usepackage{url}
\usepackage{verbatim}
\usepackage{graphicx}
\usepackage{cite}
\usepackage{amsmath,amssymb,amsthm,graphicx,booktabs}
\newtheorem{assumption}{Assumption}
\newtheorem{proposition}{Proposition}
\newtheorem{corollary}{Corollary}
\theoremstyle{remark}
\newtheorem{remark}{Remark}
\usepackage{xcolor}  
\usepackage[colorlinks=true,linkcolor=blue,urlcolor=blue,citecolor=blue]{hyperref}
\usepackage[switch]{lineno}

\hyphenation{op-tical net-works semi-conduc-tor IEEE-Xplore}
% updated with editorial comments 8/9/2021

\begin{document}
	%\linenumbers
	\title{Credible Reserve Characterization of Water-Air Coupled HVAC Systems via Work-Point Reconstruction}
	
	
	
	
	\author{Shuaihao~Jiang,
		Yanghao~Li,
		Huihui~Song,
		and Yanbin~Qu%
		\thanks{Manuscript received Month xx, 20xx; revised Month xx, 20xx; accepted Month xx, 20xx. This work was supported by the National Natural Science Foundation of China under Grant 52477085 and by the Natural Science Foundation of Shandong Province under Grant ZR2023ME162. (\textit{Corresponding author: Huihui Song})}%
		\thanks{The authors are with the School of New Energy, Harbin Institute of Technology (Weihai), Weihai 264209, Shandong, China (e-mail: shuaihaojiang@163.com; 18179477139@163.com; songhh@hitwh.edu.cn; quyanbin@hit.edu.cn).}%
	}
	
	
	% The paper headers
	% \markboth{Journal of \LaTeX\ Class Files,~Vol.~14, No.~8, August~2026}%{Shell \MakeLowercase{\textit{et al.}}: A Sample Article Using IEEEtran.cls for IEEE Journals}
	
	% \IEEEpubid{0000--0000/00\$00.00~\copyright~2021 IEEE}
	% Remember, if you use this you must call \IEEEpubidadjcol in the second
	% column for its text to clear the IEEEpubid mark.
	
	\maketitle
	
	\begin{abstract}
		Heating, ventilation, and air-conditioning (HVAC) systems in commercial buildings are increasingly expected to provide fast frequency-regulation reserves. Existing HVAC reserve-characterization methods, however, typically describe reserve around a pre-scheduled fan-side baseline and thus offer limited support for downstream aggregation and allocation. This paper proposes a work-point reconstruction framework for credible reserve characterization of water-air coupled HVAC systems. The key insight is that the water-side setpoint can reshape the air-loop operating point through water-air coupling at the coil, thereby changing the fan-power baseline and the feasible reserve range. Building on this mechanism, the proposed framework integrates a risk-calibrated thermal interface, reserve-conditioned work-point reconstruction, and deployment-oriented credibility screening to construct hourly credible reserve-cost supply functions. The results reveal that water-side coordination does not simply enlarge reserve capacity; rather, it reshapes the feasible reserve domain by relocating the fan-side baseline and changing the associated holding cost. The results also show that the practically relevant reserve region is the credible domain, since part of the scheduling-feasible reserve may become thermally fragile under deployment uncertainty. The proposed method moves HVAC reserve assessment from deterministic capacity estimation toward an hourly credible reserve-cost supply description, providing a basis for reserve quotation, bid support, and portfolio-level aggregation of building flexibility resources.
	\end{abstract}
	
	\begin{IEEEkeywords}
		Frequency regulation, building flexibility, water-air coupled HVAC, credible reserve, cost curve.
	\end{IEEEkeywords}
	
	\section{Introduction}
	\label{sec:introduction}
	
	\IEEEPARstart{T}{he} increasing penetration of converter-interfaced renewable generation has intensified the need for fast and reliable frequency-regulation resources in modern power systems\cite{ratnam2020future}. As conventional synchronous units are gradually displaced, system inertia and governor-based frequency support decline, making real-time balancing more dependent on heterogeneous regulating resources distributed across the supply, storage, and demand sides\cite{anderson2021frequency,he2024analysis}. This growing need for flexible frequency-support capability has, in  turn, accelerated the development of ancillary-service markets toward  wider resource participation and more performance-oriented procurement  frameworks\cite{rancilio2022ancillary}. Therefore, a key challenge for regulating resources is how their reserve capability can be described in a market-usable manner\cite{koltsaklis2023assessing}.
	
	Buildings are increasingly viewed as grid-interactive flexible resources because they account for a substantial portion of electricity consumption and contain controllable thermal loads with inherent energy storage\cite{tang2021flexibility}. Among building end uses, HVAC systems are particularly attractive for frequency regulation because their electrical power can be modulated through supervisory control while the building envelope and internal thermal mass buffer short-term comfort impacts\cite{hwang2023demand,ding2024flexibility}. Early studies verified that commercial-building HVAC systems can track regulation signals through fan-speed modulation while preserving indoor comfort\cite{lin2015experimental,beil2016frequency}. More recent studies have developed hierarchical continuous-control strategies\cite{wang2021hierarchical}, airflow-based virtual-battery models\cite{wang2020operating}, and frequency-regulation controllers for practical HVAC systems\cite{zhou2024hybrid}. Building on this, research on HVAC systems for frequency regulation and reserve provision has progressed from feasibility demonstration to structured reserve characterization. 
	
	A common reserve-characterization paradigm represents HVAC reserve as the admissible power deviation around a scheduled operating point, which aligns naturally with upward and downward reserve offers in ancillary-service markets such as PJM \cite{pjm2021manual11}. Vrettos et al. proposed a three-level hierarchical framework for secondary frequency reserve provision by aggregations of commercial buildings \cite{vrettos2015scheduling}, which was subsequently validated experimentally in \cite{vrettos2016experimental, vrettos2016experimental2}, confirming reliable reserve tracking and a priori capacity declaration in a real building test facility. In \cite{liu2017evaluation}, a two-step scheduling algorithm was developed that first determines an energy baseline and then quantifies regulation capacity by constraining comfort compliance at a specified percentile of the regulation signal distribution. Market-oriented studies further optimized commercial-building or HVAC participation in energy and ancillary-service markets by considering energy cost, reserve revenue, comfort impact, and expected intra-hour deployment cost\cite{aslam2023optimal}. These works established the technical foundation for HVAC regulation provision, and their reserve descriptions are predominantly baseline-centered: a nominal HVAC operating point is scheduled first, and reserve is then represented as an admissible power deviation around that point, often through fan-side modulation.
	
	Such representations are sufficient for reporting admissible capacity, but they provide limited support for downstream tasks such as aggregation, allocation, and economic valuation. For building aggregators or market operators, a reserve offer should indicate not only the admissible power range, but also whether this range remains deliverable under uncertainty\cite{liu2024continuous} and what incremental holding cost is incurred by maintaining different reserve levels \cite{chen2024aggregating}. For thermally constrained resources, flexibility provision is an additional grid service and should not compromise the primary thermal service; therefore, geometric power headroom must be further screened against thermal uncertainty, comfort constraints, and service-quality requirements\cite{maitanova2024analytical}. This perspective is consistent with recent non-residential building flexibility studies, where flexibility characterization has increasingly emphasized quantification functions, aggregation methods, and grid-integrated resource descriptions rather than isolated load-reduction or capacity values\cite{zhou2025demand}. Therefore, a market-facing HVAC reserve description should move from “how much reserve can be declared” toward “how much reserve can be credibly delivered at what cost.” 
	
	\begin{figure}[t]
		\centering
		\includegraphics[width=0.8\linewidth]{fig1_schematic.png}
		\caption{Schematic of the water--air coupled HVAC system.}
		\label{fig:hvac_system}
	\end{figure}
	
	This requirement becomes more challenging in medium- and large-scale buildings served by water–air coupled HVAC systems. In such systems, as illustrated in Fig.~\ref{fig:hvac_system}, the central heat sources or chillers supply conditioned water to air-handling units, where water-to-air coils transfer thermal energy to the supply air. In \cite{wang2022control}, studies on complete HVAC participation in energy and regulation markets have shown that air-loop-only representations miss important interactions between the heat source and the user side and degrade scheduling performance. In other words, reserve-relevant HVAC operation is rarely a purely air-side problem in large commercial buildings. At the coil, the supply-water temperature changes the thermal driving force, so the same zone thermal demand can be met by different combinations of supply-water temperature and airflow. Since fan power is a nonlinear function of airflow, these combinations correspond to different fan-power baselines and different symmetric reserve headrooms\cite{wang2019development}. Thus, the reserve-relevant fan baseline should not be treated as a fixed point; it belongs to a family of water-side-shaped operating points. This observation changes the reserve-characterization problem from scanning reserve around a fixed fan baseline to reconstructing a reserve-supporting water–air work point for each target reserve level. Therefore, the key unresolved issue is not whether HVAC systems can modulate power, but how a water-side-shaped family of fan baselines should be converted into credible reserve–cost descriptions for market use. 
	
	Motivated by this gap, this paper develops a work-point reconstruction framework for credible reserve characterization of water--air coupled HVAC systems. Instead of evaluating reserve as a power deviation around a fixed fan-side baseline, the proposed framework reconstructs the coordinated water--air operating point for each candidate reserve level. This allows HVAC reserve capability to be described jointly by the water-side-shaped fan baseline, scheduling-feasible reserve domain, deployment credibility, and incremental holding cost. The resulting hourly credible reserve--cost supply description provides a basis for reserve quotation, bid preparation, and portfolio-level aggregation of building flexibility resources.
	
	The main contributions of this paper are summarized as follows.
	\begin{enumerate}
		\item \textbf{Water-side baseline-shaping mechanism for HVAC reserve provision:} 
		This paper formalizes the supply-water temperature as a reserve-shaping variable in water--air coupled HVAC systems. It shows that, under a given thermal load, the water-side setpoint changes the coil heat-transfer condition and the required airflow, thereby shifting the fan-power baseline that determines symmetric reserve headroom.
		
		\item \textbf{Reserve-conditioned work-point reconstruction with credible-domain screening:} 
		A reserve-conditioned optimization framework is developed to reconstruct the coordinated water--air operating point for each target reserve level. The schedulable reserve domain is further screened by a risk-calibrated thermal interface and a deployment-oriented local assessment to distinguish scheduling feasibility from credible deliverability.
		
		\item \textbf{Hourly credible reserve--cost supply function:} 
		The credible reserve domain is mapped into an incremental holding-cost curve, yielding an hourly supply function that jointly reports reserve quantity, deployment credibility, and holding cost. This representation supports reserve quotation, bid preparation, and portfolio-level aggregation of building flexibility resources.
	\end{enumerate}
	
	The remainder of this paper is organized as follows. Section~\ref{sec:system_modeling} presents the reduced-order water--air coupled HVAC model and the reserve-relevant baseline shaping mechanism. Section~\ref{sec:uncertainty} develops the risk-calibrated thermal feasibility framework for credible reserve quantification. Section~\ref{sec:hourly_reserve} presents the work-point-dependent reserve supply characterization, including credibility assessment and cost-curve construction. Section~\ref{sec:Case_Study} presents the numerical results, with emphasis on the water-side baseline shaping mechanism, the credible reserve interval, and the structure of the resulting cost curves. Section~\ref{sec:conclusion} concludes this paper.
	
	
	\section{Water-Air Coupled HVAC System Model}
	\label{sec:system_modeling}
	To characterize the reserve capability of a water-air coupled HVAC system, it is necessary to describe how the fan-side baseline is formed under coordinated water- and air-side operation. In conventional air-side reserve analysis, this baseline is usually treated as a fixed scheduled operating point around which fan power is modulated \cite{vrettos2015scheduling,liu2017evaluation,aslam2023optimal}. In a water-air coupled system, however, the supply-water temperature changes the coil heat-transfer condition and thereby reshapes the airflow and fan power required to satisfy the same thermal demand. This section introduces a compact supervisory model that links indoor-temperature evolution, coil heat delivery, and electrical power.
	\subsection{Supervisory Thermal Dynamics}
	
	To evaluate whether a candidate reserve-supporting work point is thermally feasible, the indoor-temperature response over the scheduling horizon must be predicted. We use a low-order discrete-time model that maps supervisory water--air inputs and exogenous thermal disturbances to a representative controlled-zone temperature. The model serves as an optimization-oriented prediction interface, with lower-level actuator dynamics and local feedback effects absorbed into the identified dynamics and residual errors.
	
	Let \(k\) denote the basic supervisory scheduling step. We use a state-space representation of the building thermal dynamics as follows\cite{rousseau2025uncertainty}:
	\begin{equation}
		\left\{
		\begin{aligned}
			\mathbf{x}_{k+1} &= \mathbf{A}\mathbf{x}_{k}
			+ \mathbf{B}_{u}\mathbf{u}_{k}
			+ \mathbf{B}_{d}\mathbf{d}_{k}
			+ \mathbf{w}_{k}, \\
			T_{i,k} &= \mathbf{C}\mathbf{x}_{k}
			+ \mathbf{D}_{u}\mathbf{u}_{k}
			+ \mathbf{D}_{d}\mathbf{d}_{k}
			+ v_{k},
		\end{aligned}
		\right.
		\label{eq:thermal_model}
	\end{equation}
	where \(\mathbf{x}_{k}\in\mathbb{R}^{n_x}\) is the equivalent thermal state vector and \(T_{i,k}\) is the representative indoor air temperature. The supervisory control vector is
	\(\mathbf{u}_{k}=[T_{s,k},\dot m_{a,k}]^{\top}\), where \(T_{s,k}\) denotes the hydronic supply-water temperature setpoint and \(\dot m_{a,k}\) denotes the supply-air mass flow rate. The disturbance vector \(\mathbf{d}_{k}\) collects the main exogenous thermal drivers, including outdoor air temperature, solar irradiance, and internal heat gains. The terms \(\mathbf{w}_k\) and \(v_k\) denote process and measurement errors, respectively. The system matrices \(\mathbf{A}\), \(\mathbf{B}_{u}\), \(\mathbf{B}_{d}\), \(\mathbf{C}\), \(\mathbf{D}_{u}\), and \(\mathbf{D}_{d}\) can be identified from historical operating data, for example using subspace identification methods \cite{overschee1996subspace}.
	
	The noise-free propagation of \eqref{eq:thermal_model} is used in the subsequent reserve analysis to generate the nominal indoor-temperature trajectory for risk-calibrated comfort tightening. The residual uncertainty around this trajectory is quantified in Section~III, while the coil heat-delivery relation introduced next provides the physical link from supervisory inputs to baseline power.
	
	\subsection{Water--Air Coupled Baseline Formation}
	
	The supervisory thermal model in \eqref{eq:thermal_model} predicts the temperature response for a given input trajectory. To characterize reserve capability, the same water-air operating point must also be mapped to heat delivery, fan power, and source-side power.
	
	The water and air loops are coupled through the coil. As a modeling basis, the coil heat delivery is represented by a standard $\varepsilon$--NTU relation \cite{sekulic2023fundamentals}, which has also been adopted in control-oriented building HVAC modeling \cite{mantovani2014temperature}
	\begin{equation}
		Q_{h,k}
		=
		\varepsilon_k C_{\min,k}
		\left(T_{s,k}-T_{a,\mathrm{in},k}\right),
		\label{eq:ntu_heat_delivery}
	\end{equation}
	where \(Q_{h,k}\) denotes the heat delivered to the supply air, \(T_{a,\mathrm{in},k}\) is the coil inlet air temperature determined by the return-air and outdoor-air mixing condition, \(\varepsilon_k\) is the coil effectiveness, and \(C_{\min,k}\) is the smaller heat-capacity rate of the air and water sides. With \(C_{a,k}=\dot m_{a,k}c_{p,a}\) and \(C_w=\dot m_w c_{p,w}\) denoting the air- and water-side heat-capacity rates, respectively,
	\begin{equation}
		C_{\min,k}=\min\{C_{a,k},C_w\}.
		\label{eq:capacity_rates}
	\end{equation}
	The effectiveness \(\varepsilon_k\) is determined by the usual NTU and heat-capacity-ratio relation, and is therefore affected by the airflow-dependent heat-transfer condition.
	
	For the supervisory reserve analysis considered in this paper, the water-side flow rate is regulated at a lower level and treated as approximately fixed within each scheduling interval. In the studied heating operation, \(C_w\) remains larger than \(C_{a,k}\) over the admissible airflow range, so the air side becomes the limiting side of the coil heat transfer.Thus, \eqref{eq:ntu_heat_delivery} is written in the compact
	control-oriented form
	\begin{equation}
		Q_{h,k}
		=
		\varepsilon(\dot m_{a,k})
		c_{p,a}\dot m_{a,k}
		\left(T_{s,k}-T_{a,\mathrm{in},k}\right).
		\label{eq:control_heat_delivery}
	\end{equation}
	This relation shows that \(T_{s,k}\) changes the coil thermal driving force, while \(\dot m_{a,k}\) changes both the air-side heat-capacity rate and the airflow-dependent heat-transfer intensity.
	
	The same water--air operating point also determines the baseline electrical power. The total HVAC power is decomposed as
	\begin{equation}
		P_{\mathrm{tot},k}
		=
		P_{f,k}+P_{\mathrm{src},k},
		\label{eq:total_power}
	\end{equation}
	where \(P_{f,k}\) is the supply-fan power and \(P_{\mathrm{src},k}\) is the source-side power. The fan power is modeled as a fitted cubic polynomial of airflow over
	the admissible operating range,
	\begin{equation}
		P_{f,k}=g_f(\dot m_{a,k})
		=
		\sum_{i=0}^{3} a_{f,i}\dot m_{a,k}^{\,i},
		\label{eq:fan_power}
	\end{equation}
	where \(a_{f,i}\) are fitted fan coefficients. This fan-side power is the fast electrical component used for reserve deployment.
	
	The source-side power is represented as
	\begin{equation}
		P_{\mathrm{src},k}
		=
		\frac{Q_{h,k}}{\mathrm{COP}_k}
		+
		P_{\mathrm{aux}},
		\label{eq:source_power}
	\end{equation}
	where \(P_{\mathrm{aux}}\) denotes auxiliary hydronic power and \(\mathrm{COP}_k\) is the time-varying effective coefficient of performance of the heat source. In the implementation, \(\mathrm{COP}_k\) can be evaluated from a fitted performance map using the operating and weather conditions. 
	
	The reserve relevance of this coupled model can be seen by considering the heat-delivery requirement \(Q_{h,k}^{\star}\) associated with a given thermal-service condition. Any steady operating point that satisfies this requirement must fulfill
	\begin{equation}
		Q_{h,k}^{\star}
		=
		\varepsilon(\dot m_{a,k})
		c_{p,a}\dot m_{a,k}
		\left(T_{s,k}-T_{a,\mathrm{in},k}\right).
		\label{eq:equal_heat_requirement}
	\end{equation}
	Over the admissible operating range, this relation gives the airflow required under a given supply-water temperature,
	\begin{equation}
		\dot m_{a,k}^{\star}
		=
		\psi_k
		\left(
		T_{s,k};
		Q_{h,k}^{\star},T_{a,\mathrm{in},k}
		\right).
		\label{eq:airflow_family}
	\end{equation}
	Consequently, the fan-side baseline is not a fixed scalar point but a water-side-parameterized quantity,
	\begin{equation}
		P_{f,k}^{\star}(T_{s,k})
		=
		g_f
		\left(
		\psi_k
		\left(
		T_{s,k};Q_{h,k}^{\star},T_{a,\mathrm{in},k}
		\right)
		\right).
		\label{eq:fan_baseline_family}
	\end{equation}
	
	With the admissible fan-power range \([P_f^{\min},P_f^{\max}]\) induced by the airflow limits, the symmetric physical reserve around this baseline is
	\begin{equation}
		r_{k}^{\mathrm{phys}}(T_{s,k})
		=
		\min
		\left\{
		\begin{aligned}
			& P_{f,k}^{\star}(T_{s,k})-P_f^{\min},\\
			& P_f^{\max}-P_{f,k}^{\star}(T_{s,k})
		\end{aligned}
		\right\}.
		\label{eq:physical_reserve}
	\end{equation}
	Therefore, the supply-water temperature does not provide fast reserve directly; it reshapes the fan-side operating point from which fast reserve is deployed. Together with the source-side efficiency effect in \eqref{eq:source_power}, this mechanism creates the physical--economic trade-off underlying the reserve-conditioned work-point reconstruction in Section~\ref{sec:hourly_reserve}.
	
	\section{Risk-Calibrated Thermal Feasibility Interface}
	\label{sec:uncertainty}
	
	The modeling framework in Section~II provides the nominal thermal response under scheduled supervisory inputs and anticipated disturbances. In practice, the realized indoor temperature may deviate from this prediction because of weather forecast errors, unmeasured internal gains, and model--plant mismatch. This section uses residual-based conformalized quantile regression (CQR) \cite{Romano2019CQR} to quantify step-dependent temperature uncertainty and converts the resulting calibrated bounds into risk-tightened comfort constraints for reserve-feasibility assessment.
	
	\subsection{Calibrated Temperature Uncertainty Bound}
	
	Let \(t\) denote the scheduling origin, \(H\) the prediction horizon, and \(\mathcal{H}=\{1,\ldots,H\}\) the set of look-ahead steps. Given the estimated state $\hat x_t$, the planned supervisory input trajectory $\hat u_{t:t+H|t}$, and the disturbance forecast $\hat d_{t:t+H|t}$, the nominal thermal
	trajectory for \(j\in\mathcal H\) is obtained by noise-free propagation of \eqref{eq:thermal_model}:
	\begin{equation}
		\left\{
		\begin{aligned}
			\hat{\mathbf{x}}_{t+j|t}
			&=
			\mathbf{A}\hat{\mathbf{x}}_{t+j-1|t}
			+
			\mathbf{B}_{u}\hat{\mathbf{u}}_{t+j-1|t}
			+
			\mathbf{B}_{d}\hat{\mathbf{d}}_{t+j-1|t},\\
			\hat T^{\rm nom}_{i,t+j|t}
			&=
			\mathbf{C}\hat{\mathbf{x}}_{t+j|t}
			+
			\mathbf{D}_{u}\hat{\mathbf{u}}_{t+j|t}
			+
			\mathbf{D}_{d}\hat{\mathbf{d}}_{t+j|t}.
		\end{aligned}
		\right.
		\label{eq:nominal_prediction}
	\end{equation}
	Here, \((\cdot)_{t+j|t}\) denotes a prediction made at the scheduling
	origin \(t\) for the future time \(t+j\). The trajectory
	\(\hat T^{\rm nom}_{i,t+j|t}\) is used as the nominal thermal reference
	in the subsequent reserve-feasibility assessment.
	
	The remaining prediction uncertainty is modeled through the
	\(j\)-step-ahead temperature residual
	\begin{equation}
		e_{t,j}
		=
		T_{i,t+j}
		-
		\hat T^{\rm nom}_{i,t+j|t},
		\quad j\in\mathcal H .
		\label{eq:temperature_residual}
	\end{equation}
	This residual formulation separates the deterministic thermal response
	captured by the supervisory model from the uncertainty caused by
	forecast errors, unmeasured gains, and model mismatch. For each
	look-ahead step \(j\), a feature vector \(\phi_{t,j}\) is constructed
	from the information available at time \(t\), including recent
	measurements, planned inputs, disturbance forecasts, the nominal
	temperature prediction, and temporal indicators.
	
	For a target stepwise confidence level \(\alpha\in(0,1)\), the lower
	and upper residual quantile levels are set to
	\(\tau_l=(1-\alpha)/2\) and \(\tau_u=(1+\alpha)/2\). For each
	look-ahead step \(j\), two residual quantile models,
	\(\hat q_{\tau_l}^{(j)}(\phi_{t,j})\) and
	\(\hat q_{\tau_u}^{(j)}(\phi_{t,j})\), are trained offline. These
	models capture the operating-condition-dependent spread of the
	multi-step residual, allowing the uncertainty bound to vary with the
	thermal state, planned control actions, and disturbance forecasts.
	Since the raw quantile interval may not achieve the desired empirical
	coverage, a split conformal calibration step is further applied.
	
	For a held-out calibration set
	$\mathcal D_{\rm cal}^{(j)}=\{(\phi_m^{(j)},e_m^{(j)})\}_{m=1}^{n_{\rm cal}}$,
	define the nonconformity score as
	\begin{equation}
		S_m^{(j)}
		=
		\max\left\{
		\hat q_{\tau_l}^{(j)}(\phi_m^{(j)})-e_m^{(j)},
		\,
		e_m^{(j)}-\hat q_{\tau_u}^{(j)}(\phi_m^{(j)})
		\right\}.
		\label{eq:nonconformity_score}
	\end{equation}
	
	Let
	\(S_{(1)}^{(j)}\le \cdots \le S_{(n_{\rm cal})}^{(j)}\)
	be the ordered scores. The conformal inflation term at confidence
	level \(\alpha\) is defined as
	\begin{equation}
		\Delta_j(\alpha)
		=
		S_{(k_\alpha)}^{(j)},
		\qquad
		k_\alpha
		=
		\min\left\{
		\left\lceil (n_{\rm cal}+1)\alpha \right\rceil,
		n_{\rm cal}
		\right\}.
		\label{eq:conformal_inflation}
	\end{equation}
	
	The calibrated temperature bounds are then given by
	
	\begin{equation}
		\begin{aligned}
			\underline T^\alpha_{i,t+j|t}
			&=
			\hat T^{\rm nom}_{i,t+j|t}
			+
			\hat q_{\tau_l}^{(j)}(\phi_{t,j})
			-
			\Delta_j(\alpha),\\
			\overline T^\alpha_{i,t+j|t}
			&=
			\hat T^{\rm nom}_{i,t+j|t}
			+
			\hat q_{\tau_u}^{(j)}(\phi_{t,j})
			+
			\Delta_j(\alpha),
		\end{aligned}
		\quad j\in\mathcal H .
		\label{eq:calibrated_temperature_bounds}
	\end{equation}
	
	For each fixed look-ahead step \(j\), the interval
	\([\underline T^\alpha_{i,t+j|t},
	\overline T^\alpha_{i,t+j|t}]\) is interpreted as a stepwise marginal
	uncertainty bound, i.e.,
	\begin{equation}
		\Pr\left\{
		T_{i,t+j}
		\in
		\left[
		\underline T^\alpha_{i,t+j|t},
		\overline T^\alpha_{i,t+j|t}
		\right]
		\right\}
		\gtrsim
		\alpha .
		\label{eq:stepwise_coverage}
	\end{equation}
	
	Because calibration is performed separately for different look-ahead steps, this is a stepwise marginal guarantee rather than a joint trajectory-level confidence tube.
	
	\subsection{Risk-Tightened Comfort Constraint}
	
	For reserve-oriented scheduling, the calibrated temperature interval is
	converted into a deterministic tightening of the original comfort band.
	This conversion is useful because the subsequent reserve-conditioned
	optimization can then enforce thermal risk through deterministic
	constraints on the nominal trajectory, without explicitly optimizing
	over random temperature realizations.
	
	Define the lower and upper safety margins as
	\begin{equation}
		\begin{aligned}
			s^-_{t,j}(\alpha)
			&=
			-\hat q_{\tau_l}^{(j)}(\phi_{t,j})
			+
			\Delta_j(\alpha),\\
			s^+_{t,j}(\alpha)
			&=
			\hat q_{\tau_u}^{(j)}(\phi_{t,j})
			+
			\Delta_j(\alpha),
		\end{aligned}
		\quad j\in\mathcal H .
		\label{eq:safety_margins}
	\end{equation}
	The two margins may differ, thereby preserving the asymmetric and
	time-varying residual uncertainty learned from data. A larger
	\(s^-_{t,j}\) requires more distance from the lower comfort limit,
	whereas a larger \(s^+_{t,j}\) requires more distance from the upper
	comfort limit.
	
	Let \([T_i^{\min},T_i^{\max}]\) denote the original occupant comfort
	band. The corresponding risk-tightened comfort bounds for the nominal
	temperature trajectory are
	\begin{equation}
		\begin{aligned}
			\underline{\theta}^{\rm risk}_{t+j|t}
			&=
			T_i^{\min}
			+
			s^-_{t,j}(\alpha),\\
			\overline{\theta}^{\rm risk}_{t+j|t}
			&=
			T_i^{\max}
			-
			s^+_{t,j}(\alpha),
		\end{aligned}
		\quad j\in\mathcal H .
		\label{eq:risk_bounds}
	\end{equation}
	A candidate supervisory trajectory is thermally admissible at the
	stepwise confidence level \(\alpha\) if
	\begin{equation}
		\underline{\theta}^{\rm risk}_{t+j|t}
		\le
		\hat T^{\rm nom}_{i,t+j|t}
		\le
		\overline{\theta}^{\rm risk}_{t+j|t},
		\quad \forall j\in\mathcal H .
		\label{eq:risk_tightened_constraint}
	\end{equation}
	
	In each reserve-feasibility assessment, the margins in
	\eqref{eq:safety_margins} are evaluated from the offline-trained
	quantile models and conformal inflation terms and treated as fixed
	tightening coefficients. Therefore, the thermal-risk module acts as a
	calibrated outer interface to the scheduling problem rather than as an
	additional stochastic decision layer. The constraints in
	\eqref{eq:risk_tightened_constraint} characterize
	supervisory-timescale thermal feasibility, while the additional
	intra-hour temperature excursions caused by fast reserve deployment
	are screened separately in Section~IV.
	
	\section{Hourly Credible Reserve Supply Construction}
	\label{sec:hourly_reserve}
	
	Sections~II and III have established the water--air coupled supervisory model and the risk-calibrated thermal feasibility interface. This section uses these two components to construct an hourly credible reserve--cost supply function. For each target hour \(h\), the construction first obtains a risk-calibrated baseline trajectory, then reconstructs a reserve-supporting water--air work point for each candidate reserve level, screens the resulting reserve domain against deployment-induced thermal risk, and finally maps the credible domain into an incremental holding-cost curve.
	
	% \subsection{Risk-Calibrated Baseline Scheduling and Natural Reserve}
	% \label{subsec:baseline_natural_reserve}
	
	% Let $\mathcal{K}$ denote the set of supervisory scheduling
	% steps over the day, and let $\mathcal{K}^{-}$ denote the
	% corresponding state-transition index set excluding the terminal
	% step. Let $\mathcal{K}_h \subset \mathcal{K}$ denote the set of
	% scheduling steps contained in target hour $h$.
	
	% The one-step operating cost is defined as
	% \begin{equation}
		% 	\ell_k(x_k,u_k)
		% 	=
		% 	\pi_k P_{\mathrm{tot},k}
		% 	+
		% 	\rho_T
		% 	\left(
		% 	\hat T_{i,k}^{\mathrm{nom}}-T_i^{\mathrm{set}}
		% 	\right)^2 ,
		% 	\label{eq:stage_cost}
		% \end{equation}
	% where $\pi_k$ is the electricity price, $\rho_T$ is the
	% temperature-deviation penalty coefficient, and
	% $\hat T_{i,k}^{\mathrm{nom}}$ is the nominal indoor
	% temperature predicted by the supervisory model.
	
	% The risk-calibrated baseline scheduling problem is formulated as
	% \begin{subequations}
		% 	\label{prob:P0}
		% 	\begin{align}
			% 		(\mathbf{P0})\qquad
			% 		&J_0
			% 		=
			% 		\min_{\{x_k,u_k\}_{k\in\mathcal{K}}}
			% 		\quad
			% 		\sum_{k\in\mathcal{K}}
			% 		\Delta t\,\ell_k(x_k,u_k)
			% 		\label{eq:p0_obj}
			% 		\\
			% 		\mathrm{s.t.}\qquad
			% 		&
			% 		x_{k+1}
			% 		=
			% 		A x_k+B_u u_k+B_d\hat d_k, k\in\mathcal{K}^{-},
			% 		\label{eq:p0_dyn}
			% 		\\
			% 		&
			% 		\hat T_{i,k}^{\mathrm{nom}}
			% 		=
			% 		C x_k+D_u u_k+D_d\hat d_k, k\in\mathcal{K},
			% 		\label{eq:p0_temp}
			% 		\\
			% 		&
			% 		P_{\mathrm{tot},k}
			% 		=
			% 		P_{f,k}+P_{\mathrm{src},k}, k\in\mathcal{K},
			% 		\label{eq:p0_total_power}
			% 		\\
			% 		&
			% 		P_{f,k}
			% 		=
			% 		g_f(\dot m_{a,k}),k\in\mathcal{K},
			% 		\label{eq:p0_fan_power}
			% 		\\
			% 		&
			% 		P_{\mathrm{src},k}
			% 		=
			% 		g_{\mathrm{src}}
			% 		(T_{s,k},\dot m_{a,k},x_k,\hat d_k), k\in\mathcal{K},
			% 		\label{eq:p0_src_power}
			% 		\\
			% 		&
			% 		\underline{\theta}_{k}^{\mathrm{risk}}
			% 		\le
			% 		\hat T_{i,k}^{\mathrm{nom}}
			% 		\le
			% 		\overline{\theta}_{k}^{\mathrm{risk}}, k\in\mathcal{K},
			% 		\label{eq:p0_risk_comfort}
			% 		\\
			% 		&
			% 		\{u_k\}_{k\in\mathcal{K}}\in\mathcal{U},
			% 		\label{eq:p0_control_set}
			% 		\\
			% 		&
			% 		x_1=\hat x_0 .
			% 		\label{eq:p0_initial}
			% 	\end{align}
		% \end{subequations}
	
	% Here, $g_{\mathrm{src}}(\cdot)$ denotes the source-side power
	% mapping induced by the coil heat-transfer relation, the COP
	% model, and the hydronic auxiliary power introduced in
	% Section II. The set $\mathcal{U}$ collects the admissible
	% supply-water temperature and airflow limits, as well as their
	% ramping constraints.
	
	% Let $\{x_k^0,u_k^0\}_{k\in\mathcal{K}}$ be the optimal solution of $(\mathbf{P0})$, and let $P_{f,k}^0$ be the corresponding baseline fan power. The natural symmetric reserve at hour $h$
	% is defined as
	% \begin{equation}
		% 	R_h^{\mathrm{nat}}
		% 	=
		% 	\min_{k\in\mathcal{K}_h}
		% 	\min
		% 	\left\{
		% 	P_{f,k}^{0}-P_f^{\min},
		% 	\,
		% 	P_f^{\max}-P_{f,k}^{0}
		% 	\right\}.
		% 	\label{eq:natural_reserve}
		% \end{equation}
	% This quantity measures the geometric fan-power headroom
	% available around the risk-calibrated baseline trajectory. It is
	% used only as a reference, since it neither accounts for
	% reserve-induced work-point migration nor guarantees credible
	% delivery during fast intra-hour deployment.
	
	\subsection{Risk-Calibrated Baseline Scheduling and Natural Reserve}
	\label{subsec:baseline_natural_reserve}
	
	Let \(\mathcal{K}\) denote the set of supervisory scheduling steps over
	the day, and let \(\mathcal{K}^{-}\) be the corresponding
	state-transition index set excluding the terminal step. For a target
	hour \(h\), \(\mathcal{K}_h\subset\mathcal{K}\) collects the scheduling
	steps contained in that hour.
	
	The one-step operating cost is defined as
	\begin{equation}
		\ell_k(x_k,u_k)
		=
		\pi_k P_{\mathrm{tot},k}
		+
		\rho_T
		\left(
		\hat T_{i,k}^{\mathrm{nom}}-T_i^{\mathrm{set}}
		\right)^2 ,
		\label{eq:stage_cost}
	\end{equation}
	where \(\pi_k\) is the electricity price, \(\rho_T\) weights the nominal temperature deviation from the preferred setpoint.
	
	The baseline problem schedules the supervisory input trajectory, which is formulated as
	\begin{subequations}
		\label{prob:P0}
		\begin{align}
			(\mathrm{P0})\qquad
			&J_0=
			\min_{\{x_k,u_k\}_{k\in\mathcal{K}}}
			\sum_{k\in\mathcal{K}}
			\Delta t\,\ell_k(x_k,u_k)
			\label{eq:p0_obj}
			\\
			\mathrm{s.t.}\qquad
			&
			x_{k+1}
			=
			A x_k+B_u u_k+B_d\hat d_k,
			\quad k\in\mathcal{K}^{-},
			\label{eq:p0_dyn}
			\\
			&
			\hat T_{i,k}^{\mathrm{nom}}
			=
			C x_k+D_u u_k+D_d\hat d_k,
			\quad k\in\mathcal{K},
			\label{eq:p0_temp}
			\\
			&
			P_{\mathrm{tot},k}
			=
			P_{f,k}+P_{\mathrm{src},k},
			\quad k\in\mathcal{K},
			\label{eq:p0_total_power}
			\\
			&
			P_{f,k}
			=
			g_f(\dot m_{a,k}),
			\quad k\in\mathcal{K},
			\label{eq:p0_fan_power}
			\\
			&
			P_{\mathrm{src},k}
			=
			g_{\mathrm{src}}
			(T_{s,k},\dot m_{a,k},x_k,\hat d_k),
			\quad k\in\mathcal{K},
			\label{eq:p0_src_power}
			\\
			&
			\theta_{k}^{\mathrm{risk},-}
			\le
			\hat T_{i,k}^{\mathrm{nom}}
			\le
			\theta_{k}^{\mathrm{risk},+},
			\quad k\in\mathcal{K},
			\label{eq:p0_risk_comfort}
			\\
			&
			\{u_k\}_{k\in\mathcal{K}}\in\mathcal{U},
			\label{eq:p0_control_set}
			\\
			&
			x_1=\hat x_0 .
			\label{eq:p0_initial}
		\end{align}
	\end{subequations}
	
	Here, \(g_{\mathrm{src}}(\cdot)\) denotes the source-side power mapping
	induced by the coil heat-transfer relation, the heat-source performance
	map, and auxiliary hydronic power. The set \(\mathcal{U}\) is the
	trajectory-level admissible set for the supervisory controls, including
	the supply-water temperature and airflow limits and their ramping
	constraints.
	
	Let \(\{x_k^0,u_k^0\}_{k\in\mathcal{K}}\) be the optimal solution of
	\((\mathrm{P0})\), and let \(P_{f,k}^0\) be the corresponding baseline
	fan power. The natural symmetric reserve around this baseline at hour
	\(h\) is defined as
	\begin{equation}
		R_h^{\mathrm{nat}}
		=
		\min_{k\in\mathcal{K}_h}
		\min
		\left\{
		P_{f,k}^{0}-P_f^{\min},
		\,
		P_f^{\max}-P_{f,k}^{0}
		\right\}.
		\label{eq:natural_reserve}
	\end{equation}
	This quantity measures the geometric fan-power headroom available
	around the risk-calibrated baseline trajectory. It is used only as a
	baseline reference, because it neither allows reserve-induced
	work-point migration nor screens the delivery risk caused by fast
	intra-hour deployment.
	
	% \subsection{Reserve-Conditioned Work-Point Reconstruction}
	% \label{subsec:work_point_reconstruction}
	
	% For a target hour $h$ and a candidate symmetric reserve level
	% $R\ge0$, the fan must retain upward and downward power
	% headroom throughout that hour. Therefore,
	% \begin{equation}
		% 	P_f^{\min}+R
		% 	\le
		% 	P_{f,k}
		% 	\le
		% 	P_f^{\max}-R,
		% 	\qquad
		% 	k\in\mathcal{K}_h .
		% 	\label{eq:reserve_headroom}
		% \end{equation}
	% Since the fan-power mapping $P_f=g_f(\dot m_a)$ is
	% monotonically increasing in the operating range of interest,
	% \eqref{eq:reserve_headroom} equivalently restricts the feasible
	% airflow interval during the target hour.
	
	% The reserve-conditioned work-point reconstruction problem is
	% then formulated as
	% \begin{subequations}
		% 	\label{prob:P1}
		% 	\begin{align}
			% 		(\mathbf{P1}(R_h))\qquad
			% 		&J_h(R)=\min_{\{x_k,u_k\}_{k\in\mathcal{K}}}
			% 		\sum_{k\in\mathcal{K}}
			% 		\Delta t\,\ell_k(x_k,u_k)
			% 		\label{eq:p1_obj}
			% 		\\
			% 		\mathrm{s.t.}\qquad
			% 		&
			% 		\eqref{eq:p0_dyn}--\eqref{eq:p0_initial},
			% 		\label{eq:p1_common}
			% 		\\
			% 		P_f^{\min}+R
			% 			&\le
			% 		P_{f,k}
			% 		\le
			% 		P_f^{\max}-R,
			% 		\qquad
			% 		k\in\mathcal{K}_h .
			% 		\label{eq:p1_headroom}
			% 	\end{align}
		% \end{subequations}
	
	% Formulation $(\mathbf{P1}_h(R))$ inherits the thermal
	% dynamics, power model, risk-tightened comfort constraints,
	% control limits, and initial condition from $(\mathbf{P0})$, and
	% adds the target-hour reserve-headroom constraint
	% \eqref{eq:p1_headroom}. This constraint does not merely check
	% whether the baseline remains feasible. Instead, it embeds the
	% candidate reserve level into the fan-power feasible set and,
	% through the water--air coupling, induces a new optimal
	% reserve-supporting work point.
	
	% Let
	% \begin{equation}
		% 	z_h^\star(R)
		% 	=
		% 	\left\{
		% 	x_k^\star(R),\,
		% 	u_k^\star(R),\,
		% 	P_{f,k}^\star(R),\,
		% 	\hat T_{i,k}^{\mathrm{nom},\star}(R)
		% 	\right\}_{k\in\mathcal{K}_h}
		% 	\label{eq:work_point}
		% \end{equation}
	% denote the target-hour reserve-conditioned work point. The
	% corresponding schedulable reserve upper bound is defined as
	% \begin{equation}
		% 	R_h^{\mathrm{sch}}
		% 	=
		% 	\sup
		% 	\left\{
		% 	R\ge0
		% 	\mid
		% 	(\mathbf{P1}_h(R))\ \text{is feasible}
		% 	\right\}.
		% 	\label{eq:schedulable_reserve}
		% \end{equation}
	% This quantity characterizes the maximum reserve that can be
	% held at the supervisory scheduling layer under risk-tightened
	% comfort constraints and actuator limits.
	
	\subsection{Reserve-Conditioned Work-Point Reconstruction}
	\label{subsec:work_point_reconstruction}
	
	For a target hour \(h\) and a candidate hourly symmetric reserve level
	\(R\ge0\), the fan must retain upward and downward power headroom at
	each supervisory step in that hour. Therefore,
	\begin{equation}
		P_f^{\min}+R
		\le
		P_{f,k}
		\le
		P_f^{\max}-R,
		\qquad
		k\in\mathcal{K}_h .
		\label{eq:reserve_headroom}
	\end{equation}
	Since \(g_f(\cdot)\) is monotonically increasing over the admissible
	airflow range, \eqref{eq:reserve_headroom} equivalently contracts the
	allowable airflow interval during the target hour.
	
	The reserve-conditioned work-point reconstruction problem is then
	formulated as
	\begin{subequations}
		\label{prob:P1}
		\begin{align}
			(\mathrm{P1}_h(R))\qquad
			J_h(R)
			&=
			\min_{\{x_k,u_k\}_{k\in\mathcal{K}}}
			\sum_{k\in\mathcal{K}}
			\Delta t\,\ell_k(x_k,u_k)
			\label{eq:p1_obj}\\
			\mathrm{s.t.}\qquad
			\eqref{eq:p0_dyn}&--\eqref{eq:p0_initial},\eqref{eq:reserve_headroom}
		\end{align}
	\end{subequations}
	
	Problem \((\mathrm{P1}_h(R))\) inherits the thermal dynamics, power
	model, risk-tightened comfort constraints, control limits, and initial
	condition from \((\mathrm{P0})\), and adds the target-hour headroom
	constraint \eqref{eq:p1_headroom}. Rather than checking headroom around
	the baseline trajectory, \eqref{eq:p1_headroom} contracts the target-hour
	fan-power feasible set and induces a re-optimized water--air work point
	for the candidate reserve level.
	
	Let
	\begin{equation}
		z_h^\star(R)
		=
		\left\{
		x_k^\star(R),\,
		u_k^\star(R),\,
		P_{f,k}^\star(R),\,
		\hat T_{i,k}^{\mathrm{nom},\star}(R)
		\right\}_{k\in\mathcal{K}_h}
		\label{eq:work_point}
	\end{equation}
	denote the target-hour reserve-conditioned work-point segment. The
	corresponding schedulable reserve upper bound is defined as
	\begin{equation}
		R_h^{\mathrm{sch}}
		=
		\sup
		\left\{
		R\ge0
		\mid
		(\mathrm{P1}_h(R))\ \text{is feasible}
		\right\}.
		\label{eq:schedulable_reserve}
	\end{equation}
	This quantity is the maximum reserve that can be held at the
	supervisory scheduling layer under risk-tightened comfort constraints
	and actuator limits. Unlike \(R_h^{\mathrm{nat}}\), which is evaluated
	around the baseline trajectory, \(R_h^{\mathrm{sch}}\) allows the
	water--air trajectory to be re-optimized for the reserve requirement.
	\newpage
	\subsection{Deployment-Credible Reserve Screening}
	\label{subsec:credible_screening}
	
	A reserve level feasible in \((\mathrm{P1}_h(R))\) may still fail
	to be credibly deliverable during fast regulation deployment.
	To screen such deployment-induced risk, a local surrogate is
	constructed around \(z_h^\star(R)\).
	
	Let \(\tau\in\mathcal{M}_h\) denote the fast-time index within target hour \(h\), where \(\mathcal{M}_h\) is the set of discrete substeps during the target hour. Let \(a_\tau \in [-1, 1]\) be the normalized regulation signal, and let \(k(\tau) \in \mathcal{K}_h\) denote the supervisory scheduling step corresponding to fast-time instant \(\tau\). The fan-power reference during deployment is written as
	\begin{equation}
		P_{f,\tau}^{\mathrm{ref}}
		=
		P_{f,k(\tau)}^\star(R)+R a_\tau .
		\label{eq:fan_ref}
	\end{equation}
	where \(a_\tau\) represents the normalized regulation signal, and positive/negative values of \(a_\tau\) indicate increasing or decreasing fan power, respectively.
	
	Linearizing the local thermal response around the reserve-conditioned work point yields the following state-space approximation:
	\begin{equation}
		\left\{
		\begin{aligned}
			\delta x_{\tau+1}
			&=
			A_{h,\tau}(R)\delta x_\tau
			+
			B_{h,\tau}(R)R a_\tau
			+
			E_{h,\tau}(R)\nu_\tau ,
			\label{eq:local_state}
			\\
			\delta T_{i,\tau}
			&=
			H_{h,\tau}(R)\delta x_\tau ,
		\end{aligned}
		\right.
	\end{equation}
	where \(\delta x_\tau\) represents the deviation from the reserve-conditioned state, \(\nu_\tau\) models high-frequency disturbances and local model mismatch, and \(H_{h,\tau}(R)\) maps state deviations to indoor temperature deviations.
	
	Let $\sigma_{a,h}^2$ denote the variance of the regulation
	signal during hour $h$, and let $\Sigma_{\nu,h}$ denote the
	local disturbance covariance. The covariance of local state
	deviations evolves as
	\begin{align}
		\Sigma_{\tau+1}
		=&\,
		A_{h,\tau}(R)\Sigma_\tau A_{h,\tau}^{\top}(R)
		+
		B_{h,\tau}(R)R^2\sigma_{a,h}^2 B_{h,\tau}^{\top}(R)
		\nonumber\\
		&+
		E_{h,\tau}(R)\Sigma_{\nu,h}E_{h,\tau}^{\top}(R).
		\label{eq:covariance_prop}
	\end{align}
	The deployment-induced thermal-risk margin is then defined as
	\begin{equation}
		d_h(R,\beta)
		=
		z_{\beta}
		\max_{\tau\in\mathcal{M}_h}
		\sqrt{
			H_{h,\tau}(R)\Sigma_\tau H_{h,\tau}^{\top}(R)
		},
		\label{eq:deployment_margin}
	\end{equation}
	where $\beta$ is the deployment-level confidence parameter.
	
	The reserve level $R$ is regarded as deployment-credible if
	\begin{equation}
		\underline{\theta}_{k}^{\mathrm{risk}}
		+
		d_h(R,\beta)
		\le
		\hat T_{i,k}^{\mathrm{nom},\star}(R)
		\le
		\overline{\theta}_{k}^{\mathrm{risk}}
		-
		d_h(R,\beta),
		\qquad
		k\in\mathcal{K}_h .
		\label{eq:credible_condition}
	\end{equation}
	Because \eqref{eq:p1_headroom} together with
	$a_\tau\in[-1,1]$ already guarantees instantaneous fan-power
	feasibility, the screening step focuses on thermal credibility.
	Airflow feasibility is checked numerically through the inverse
	fan map and omitted here for brevity.
	
	The credible reserve upper bound is defined as
	\begin{equation}
		R_h^{\mathrm{cred}}
		=
		\sup
		\left\{
		R\in[0,R_h^{\mathrm{sch}}]
		\mid
		\eqref{eq:credible_condition}\ \text{holds}
		\right\}.
		\label{eq:credible_reserve}
	\end{equation}
	Accordingly,
	\begin{equation}
		0
		\le
		R
		\le
		R_h^{\mathrm{cred}}
		\le
		R_h^{\mathrm{sch}} .
		\label{eq:reserve_hierarchy}
	\end{equation}
	The schedulable bound $R_h^{\mathrm{sch}}$ reflects
	supervisory-layer feasibility, whereas $R_h^{\mathrm{cred}}$
	further excludes reserve levels that may violate thermal
	admissibility during fast deployment.
	
	\subsection{Holding-Cost Curve and Supply Function}
	\label{subsec:cost_supply_function}
	
	For any reserve level
	$R\in[0,R_h^{\mathrm{cred}}]$, the trajectory-level holding cost
	is defined as
	\begin{equation}
		C_h(R)
		=
		J_h(R)-J_0 .
		\label{eq:holding_cost}
	\end{equation}
	Here, $J_0$ is the risk-calibrated baseline scheduling cost
	without reserve commitment, and $J_h(R)$ is the optimal cost
	when reserve level $R$ is held at target hour $h$. Hence,
	$C_h(R)$ represents the supervisory-layer opportunity cost
	induced by the reserve commitment.
	
	Since $J_h(R)$ is obtained from a day-ahead scheduling
	problem, $C_h(R)$ captures not only the target-hour work-point
	migration, but also the redistribution between fan-side and
	source-side power and the downstream thermal recovery
	associated with reserve support. The cost defined in
	\eqref{eq:holding_cost} does not include real-time mileage,
	performance scores, or market-specific settlement terms.
	
	The hourly credible reserve supply function is therefore
	defined as
	\begin{equation}
		\mathcal{S}_h^{\mathrm{cred}}
		=
		\left\{
		(R,C_h(R))
		\mid
		0\le R\le R_h^{\mathrm{cred}}
		\right\}.
		\label{eq:final_supply_function}
	\end{equation}
	When needed, a marginal offer-price reference can be obtained
	from finite differences of $C_h(R)$ over a reserve scan grid.
	
	\section{Case Study}
	\label{sec:Case_Study}
	
	\subsection{Simulation Setup and Benchmarks}
	\label{subsec:simulation_setup}
	
	The case study considers a medium-sized water--air coupled HVAC system serving a public exhibition building with a floor area of approximately \(10{,}886~\mathrm{m}^2\). The system configuration and main physical parameters are adapted from a field air-source heat-pump district heating system located in Rongcheng, Shandong Province, China, of which the detailed descreption can be found in \cite{jiang2026modeling}. Historical operation data include outdoor temperature, indoor temperature, and plant-side and supervisory-level operating measurements, with a sampling interval of 15 min. The data period spans Jan. 1--31, 2026. Internal heat gains caused by occupants, lighting, and equipment are not directly measured and are therefore estimated using a function-zone-based approach \cite{wang2006parameter}. 
	
	A second-order discrete-time state-space model is identified using the N4SID algorithm and used as the supervisory thermal model. The training, calibration, and testing datasets are split chronologically with a ratio of 60\%/20\%/20\%. The stepwise marginal coverage level of the CQR-based thermal uncertainty bound is set to \(\alpha=0.90\). For the subsequent reserve evaluation, the proposed water--air coordinated reconstruction is compared with an air-only benchmark, where the supply-water temperature follows the baseline schedule and the reserve capability is characterized only through fan-side modulation.
	
	Jan. 27, 2026 is selected as an independent winter test day. The outdoor temperature, solar irradiance, monthly RegD statistics associated with the test day, and day-ahead electricity prices are shown in Fig.~\ref{fig:test_day_profiles}\cite{PJM_MarketBasedRegulation}. All scheduling and screening problems are solved over a 24-h horizon with a 15-min discretization. The resulting nonlinear programming problems are solved in MATLAB using \texttt{fmincon} with the SQP algorithm \cite{MathWorks_OptimizationToolbox}. The main simulation settings are summarized in Table~\ref{tab:simulation_setup}.
	
	\begin{figure}[!ht]
		\centering
		\includegraphics[width=1\columnwidth]{fig3_input.pdf}  
		\caption{Actual indoor temperature trajectory and calibrated 90\% prediction interval on a representative test day.}
		\label{fig:test_day_profiles}
	\end{figure}
	
	\begin{table}[t]
		\caption{Main Simulation Settings}
		\label{tab:simulation_setup}
		\centering
		\footnotesize
		\setlength{\tabcolsep}{3pt}
		\renewcommand{\arraystretch}{1.06}
		\begin{tabular}{p{0.45\columnwidth} p{0.47\columnwidth}}
			\toprule
			\textbf{Item} & \textbf{Value} \\
			\midrule
			\multicolumn{2}{l}{\textit{Data and model settings}} \\
			Building floor area & \(10886~\mathrm{m}^2\) \\
			Data period & Jan. 1--31, 2026 \\
			Independent test day & Feb. 14, 2026 \\
			Time step / horizon & 15 min / 24 h \\
			Thermal model & Second-order N4SID model \\
			Data split & 60\%/20\%/20\% chronological split \\
			\midrule
			\multicolumn{2}{l}{\textit{Comfort, risk, and reserve settings}} \\
			Comfort band / setpoint & \([20,24]^\circ\mathrm{C}\) / \(22^\circ\mathrm{C}\) \\
			CQR coverage level & \(\alpha=0.90\) \\
			Deployment confidence level & \(\beta=0.90\) \\
			Reserve scan grid size & \(n_{\mathrm{scan}}=60\) \\
			Deployment discretization & Four 15-min substeps per hour \\
			\midrule
			\multicolumn{2}{l}{\textit{HVAC physical settings}} \\
			Supply-water temperature range & \([35,50]^\circ\mathrm{C}\) \\
			Airflow range & \([6,24]~\mathrm{kg/s}\) \\
			Ramping limits of \(T_s\) and \(\dot m_a\) & \(10^\circ\mathrm{C}\), \(10~\mathrm{kg/s}\) per step \\
			Maximum heating capacity & \(650~\mathrm{kW}\) \\
			Coil effectiveness & \(\varepsilon=0.9\) \\
			Heat-pump COP model & \(\max\{1.2,\,5.148-0.075(T_s-T_o)\}\) \\
			Fan power model & Cubic polynomial in \(\dot m_a\) \\
			Fan coefficients & \((10.133,10.088,-59.566,208.298)\) \\
			\midrule
			\multicolumn{2}{l}{\textit{Market and optimization settings}} \\
			Regulation signal & Historical PJM RegD \\
			Electricity price & PJM day-ahead price \\
			Objective weights & \(\rho_{\mathrm{temp}}=2\), \(\rho_{\mathrm{slack}}=10^4\) \\
			Optimization solver & MATLAB \texttt{fmincon}, SQP \\
			\bottomrule
		\end{tabular}
	\end{table}
	
	
	\begin{table}[t]
		\caption{Main Simulation Settings}
		\label{tab:simulation_setup}
		\centering
		\footnotesize
		\setlength{\tabcolsep}{3pt}
		\renewcommand{\arraystretch}{1.05}
		\begin{tabular}{p{0.48\columnwidth} p{0.43\columnwidth}}
			\toprule
			\textbf{Parameter} & \textbf{Value / Source} \\
			\midrule
			\multicolumn{2}{l}{\textit{Building, data, and thermal interface}} \\
			Building floor area & \(10{,}886~\mathrm{m}^2\) \\
			Data period & Jan. 1--31, 2026 \\
			Independent test day & Feb. 14, 2026 \\
			Sampling interval & 15 min \\
			Thermal prediction model & Second-order N4SID model \\
			Train/calibration/test split & \(60\%/20\%/20\%\) chronological split \\
			CQR coverage level & \(\alpha=0.90\) \\
			\midrule
			\multicolumn{2}{l}{\textit{Comfort and HVAC settings}} \\
			Comfort band and setpoint & \([20,24]^\circ\mathrm{C}\), \(T_i^{\mathrm{set}}=22^\circ\mathrm{C}\) \\
			Supply-water temperature range & \([35,50]^\circ\mathrm{C}\) \\
			Airflow range & \([6,24]~\mathrm{kg/s}\) \\
			Maximum heating capacity & \(650~\mathrm{kW}\) \\
			\midrule
			\multicolumn{2}{l}{\textit{Reserve screening and optimization}} \\
			Regulation signal statistics & Historical PJM RegD data \cite{PJM_MarketBasedRegulation} \\
			Deployment representation & Four 15-min substeps per hour \\
			Scheduling horizon and step & 24 h, 15-min step \\
			Reserve scan grid size & \(n_{\mathrm{scan}}=60\) \\
			Electricity price & PJM day-ahead price \\
			Optimization solver & MATLAB \texttt{fmincon}, SQP \cite{MathWorks_OptimizationToolbox} \\
			\bottomrule
		\end{tabular}
	\end{table}
	
	
	% \subsection{Simulation Setup}
	% \label{subsec:simulation_setup}
	
	% The simulation study considers a medium—sized water--air coupled HVAC system serving a public exhibition building with a floor area of approximately $10{,}886~\mathrm{m}^2$. The plant configuration and key physical parameters are adapted from a field air-source heat pump district heating system (ASHP-DHS) located in Rongcheng, Shandong Province, China \cite{jiang2026modeling}. Operational data, including outdoor air temperature, indoor air temperature, and plant-/supervisory-level operating measurements, were collected on-site at a 15-min resolution over a continuous 31-day period during the 2025 heating season(2026.1.1-2026.1.31). Internal heat gains from occupants, lighting, and equipment are not directly metered; instead, they are estimated using the zone-type-based method in \cite{wang2006parameter}, which derives representative weekday and weekend profiles from building function and occupancy schedules. Solar irradiance data are obtained from a local meteorological station and incorporated into the dataset as an exogenous disturbance input. Together, these measured and estimated sequences form a five-channel input vector
	% \begingroup
	% \setlength{\abovedisplayskip}{2pt}
	% \setlength{\belowdisplayskip}{2pt}
	% \setlength{\abovedisplayshortskip}{2pt}
	% \setlength{\belowdisplayshortskip}{2pt}
	% \[
	% [T_s,\ \dot{m}_a,\ T_o,\ I_{\mathrm{sol}},\ Q_{\mathrm{int}}],
	% \]
	% \endgroup
	% and a single-channel indoor temperature output, all synchronized at 15-min intervals.
	
	% Using this dataset, a second-order discrete-time state-space model is identified by the N4SID algorithm. The model captures the dominant thermal dynamics at the supervisory timescale and operates directly at the native 15-min resolution of the optimization problem formulated in Section~IV. The risk-tightened comfort bounds, denoted by $\underline{\theta}_k^{\mathrm{risk}}$ and $\bar{\theta}_k^{\mathrm{risk}}$, are obtained from the CQR-plus-conformal calibration procedure described in Section~III. The uncertainty model is trained using a chronological $60\%/20\%/20\%$ split of the full dataset and targets a stepwise marginal coverage level of $\alpha=0.90$.
	
	% The day-ahead simulations are conducted for a representative winter day (2026.2.14), using the corresponding 15-min profiles of ambient temperature, internal gains, and solar irradiance extracted from the same dataset. The indoor temperature setpoint is fixed at $22^\circ\mathrm{C}$, and the original comfort band is specified as $[20^\circ\mathrm{C},\,24^\circ\mathrm{C}]$. In the credibility-screening stage, each hour is further divided into four 15-min substeps, and the intra-hour deployment uncertainty is characterized using historical PJM RegD statistics\cite{PJM_MarketBasedRegulation}. The electricity price follows the PJM day-ahead market. The fast-timescale regulation signal $a_\tau$ used in the credibility screening is represented through the hourly variance of the PJM RegD signal. The reserve-conditioned scheduling problem in \eqref{eq:schedulable_reserve_bound} is solved in MATLAB using \texttt{fmincon} with the SQP algorithm\cite{MathWorks_OptimizationToolbox}. A uniform reserve scan is adopted over a finite set of candidate reserve levels. Key system parameters and simulation settings are summarized in Table~\ref{tab:simulation_setup}.
	
	% \begin{table}[t]
		% 	\caption{Key System Parameters and Simulation Settings}
		% 	\label{tab:simulation_setup}
		% 	\centering
		% 	\footnotesize
		% 	\setlength{\tabcolsep}{3pt}
		% 	\renewcommand{\arraystretch}{1.05}
		% 	\begin{tabular}{p{0.48\columnwidth} p{0.43\columnwidth}}
			% 		\toprule
			% 		\textbf{Parameter} & \textbf{Value / Source} \\
			% 		\midrule
			% 		\multicolumn{2}{l}{\textit{Building, data, and thermal interface}} \\
			% 		Building floor area & $10{,}886~\mathrm{m}^2$ \\
			% 		Data used for model training & 31 days, 15-min resolution \\
			% 		Representative simulation day & Feb. 14, 2026 \\
			% 		Input/output variables & $[T_s,\dot m_a,T_o,I_{\mathrm{sol}},Q_{\mathrm{int}}]\rightarrow T_i$ \\
			% 		Thermal prediction model & Second-order N4SID model \\
			% 		Train/calibration/test split & $60\%/20\%/20\%$ chronological split \\
			% 		CQR coverage level & $\beta=0.90$ \\
			% 		Comfort band and setpoint & $[20,24]^\circ\mathrm{C}$, $T_i^{\mathrm{set}}=22^\circ\mathrm{C}$ \\
			% 		\midrule
			% 		\multicolumn{2}{l}{\textit{HVAC constraints and power model}} \\
			% 		Supply-water temperature range & $[35,50]^\circ\mathrm{C}$ \\
			% 		Airflow range & $[6,24]~\mathrm{kg/s}$ \\
			% 		Fan power model & Cubic polynomial $P_f=g_f(\dot m_a)$ \\
			% 		Source-side power model & COP-based heat-pump power model \\
			% 		Maximum heating capacity & $650~\mathrm{kW}$ \\
			% 		\midrule
			% 		\multicolumn{2}{l}{\textit{Reserve screening and optimization}} \\
			% 		Regulation signal statistics & Historical PJM RegD data \\
			% 		Deployment representation & Four 15-min substeps per hour \\
			% 		Scheduling horizon and step & 24 h, 15-min step \\
			% 		Reserve scan grid size & $n_{\mathrm{scan}}=60$ \\
			% 		Price and objective components & PJM day-ahead price; energy, comfort, and slack costs \\
			% 		Optimization solver & MATLAB \texttt{fmincon}, SQP \\
			% 		\bottomrule
			% 	\end{tabular}
		% \end{table}
	
	% \begin{figure}[!ht]
		%     \centering
		%     \includegraphics[width=1\columnwidth]{fig3_input.pdf}  
		%     \caption{Actual indoor temperature trajectory and calibrated 90\% prediction interval on a representative test day.}
		%     \label{fig:testday_cqr_interval}
		% \end{figure}
	
	\subsection{Validation of the Risk-Aware Thermal Interface}
	\label{subsec:validation_risk_interface}
	
	The subsequent credibility screening relies on the quality of the risk-tightened comfort bounds, $\underline{\theta}_k^{\mathrm{risk}}$ and $\overline{\theta}_k^{\mathrm{risk}}$, generated by the CQR-plus-conformal procedure. To verify their reliability, the proposed thermal risk interface is evaluated on an independent test day that is excluded from both model training and conformal calibration. As shown in Fig.~\ref{fig:Temperature_under_90_converge}, the actual indoor temperature remains within the calibrated 90\% prediction interval throughout the representative test day. Over the entire test set, the nominal prediction achieves an MAE of $0.0742^\circ\mathrm{C}$ and an RMSE of $0.1223^\circ\mathrm{C}$, while the empirical coverage under the nominal 90\% confidence level reaches $95.94\%$. These results indicate that the reduced-order supervisory model captures the dominant thermal dynamics with satisfactory accuracy, and that the conformalized interval provides an empirically reliable uncertainty envelope on unseen data.
	
	More importantly, the risk-aware thermal interface is not a fixed conservative margin, but a time-varying and direction-dependent tightening of the comfort region. As shown in Fig.~\ref{fig:cqr_vs_gaussian_margin}, the CQR-plus-conformal margins remain close to the Gaussian margins during relatively stable periods, but deviate significantly during the morning transition and the noon-to-afternoon high-load period. In particular, the upper and lower margins evolve differently, revealing that the prediction uncertainty is both time-varying and asymmetric. This feature is critical for reserve qualification: a fixed symmetric Gaussian margin may be unnecessarily conservative during low-uncertainty periods, while underestimating one-sided thermal risk during high-uncertainty periods. By contrast, the proposed data-driven interface reshapes the admissible thermal region adaptively in both time and direction, which provides a more appropriate thermal basis for the subsequent credibility screening.
	
	\begin{figure}[!ht]
		\centering
		\includegraphics[width=0.85\columnwidth]{fig4_Temperature_under_90_converge.pdf}
		\caption{Actual indoor temperature and calibrated 90\% prediction interval on a representative test day.}
		\label{fig:Temperature_under_90_converge}
	\end{figure}
	
	\begin{figure}[!ht]
		\centering
		\includegraphics[width=0.9\columnwidth]{fig5_cqr_vs_gau.pdf}
		\caption{Comparison of Gaussian and CQR-plus-conformal safety margins on the representative test day.}
		\label{fig:cqr_vs_gaussian_margin}
	\end{figure}
	
	\subsection{Evidence of Water-Side Baseline Shaping}
	\label{subsec:baseline_shaping_evidence}
	
	To verify the baseline-shaping role of the water side, two supervisory strategies are compared under the same outdoor conditions, internal heat gains, electricity prices, and risk-tightened comfort constraints. The first is the proposed coordinated water--air strategy, in which the supply-water temperature $T_s$ and the airflow rate $\dot m_a$ are jointly optimized. The second is an air-only benchmark, in which $T_s$ is fixed at $42^\circ\mathrm{C}$ and only $\dot m_a$ is optimized. This comparison isolates the baseline-shaping effect of water-side scheduling under comparable thermal-service requirements.
	
	Fig.~\ref{fig:control_compare}(a) shows the indoor-temperature trajectories and the risk-tightened comfort interval under the two strategies. Both strategies keep the indoor temperature within the risk-tightened interval and maintain the trajectory around the $22^\circ\mathrm{C}$ setpoint. Compared with the air-only benchmark, the coordinated strategy exhibits a slightly wider but still bounded temperature fluctuation, indicating that part of the available comfort margin is used to support the water--air work-point adjustment. This provides a comparable thermal-service basis for interpreting the subsequent differences in airflow and power decomposition.
	
	Fig.~\ref{fig:control_compare}(b) further shows the corresponding baseline supervisory controls. In the air-only benchmark, $T_s$ is fixed at $42^\circ\mathrm{C}$, and the intra-day thermal-load variation is mainly accommodated by the airflow. Under the coordinated strategy, $T_s$ is used as an additional scheduling degree of freedom, allowing similar indoor-temperature trajectories to be maintained through different water--air work points. During 8--16 h, the increase in $T_s$ strengthens the coil-side thermal driving condition and reduces the airflow required for thermal balance. This behavior shows that the fan operating baseline is influenced by both the thermal load and the water-side heat-transfer condition.
	
	\begin{figure}[!t]
		\centering
		\includegraphics[width=0.9\columnwidth]{fig6_control_compare.pdf}
		\caption{Comparison of indoor temperature and baseline supervisory controls under the coordinated water--air strategy and the air-only benchmark. 
			(a) Indoor-temperature trajectories with the risk-tightened comfort interval. 
			(b) Baseline airflow and supply-water temperature trajectories; the shaded region marks the 8--16 h interval.}
		\label{fig:control_compare}
	\end{figure}
	
	The corresponding baseline power redistribution is shown in Fig.~\ref{fig:power_breakdown_compare}. Over the 24-h horizon, the coordinated water--air strategy reduces the average heat-pump power from $117.10~\mathrm{kW}$ to $95.10~\mathrm{kW}$, while increasing the average fan power from $34.74~\mathrm{kW}$ to $51.32~\mathrm{kW}$. The average total HVAC power is slightly reduced from $151.84~\mathrm{kW}$ to $146.42~\mathrm{kW}$. These results suggest that water-side participation changes the baseline allocation between source-side heat delivery and air-side transport. Under comparable indoor thermal service, the supply-water temperature modifies the airflow requirement through the coil heat-transfer condition and relocates the fan-power baseline to a different operating region.
	
	\begin{figure}[t]
		\centering
		\includegraphics[width=1\columnwidth]{fig7_Power_breakdown.pdf}
		\caption{Comparison of baseline power decomposition under the coordinated water--air strategy and the air-only benchmark.}
		\label{fig:power_breakdown_compare}
	\end{figure}
	
	Together, Fig.~\ref{fig:control_compare} and Fig.~\ref{fig:power_breakdown_compare} support the water-side baseline-shaping mechanism introduced in Section~II. The water-side setpoint modifies the coil heat-transfer condition, changes the airflow trajectory required for thermal comfort, and thereby shifts the fan-power baseline. Since the symmetric reserve headroom directly depends on the location of this fan-power baseline, the next subsection quantifies the relationship among natural reserve, scheduling-feasible reserve, and credible reserve in a unified framework.
	
	\subsection{From Geometric Feasibility to Credible Reserve}
	
	Fig.~\ref{fig:credible_reserve} compares the three reserve boundaries considered in this study: the natural reserve $R_{\mathrm{nat}}(h)$, the maximum scheduling-feasible reserve $R_{\max}(h)$, and the maximum credible reserve $R_{\mathrm{cred}}(h)$. The natural reserve represents the symmetric fan-side geometric headroom around the coordinated water--air baseline. The maximum scheduling-feasible reserve is the feasible upper bound obtained after reconstructing the water--air work point under a candidate reserve requirement, whereas the maximum credible reserve further accounts for intra-hour regulation deployment uncertainty and the risk-tightened comfort bounds. Over the 24-h horizon, work-point reconstruction increases the average reserve upper bound from $43.57~\mathrm{kW}$ to $67.05~\mathrm{kW}$, yielding an average expansion of $23.49~\mathrm{kW}$ relative to the natural reserve. After credibility screening, the average reserve is reduced to $57.13~\mathrm{kW}$, corresponding to an average contraction of $9.92~\mathrm{kW}$, or $14.8\%$ of the scheduling-feasible reserve. The contraction is more pronounced during 9--17 h, where the average reduction reaches $19.95~\mathrm{kW}$, approximately $33.0\%$ of the scheduling-feasible interval. This mid-day contraction indicates that a larger portion of the scheduling-feasible reserve is exposed to deployment-induced thermal risk, motivating a closer examination of the work-point reconstruction and credibility screening process in a representative hour.
	
	\begin{figure}[htbp]
		\centering
		\includegraphics[width=0.8\columnwidth]{fig9_credible_reserve.pdf}
		\caption{Hourly comparison of natural reserve, scheduling-feasible reserve, and credible reserve.}
		\label{fig:credible_reserve}
	\end{figure}
	
	Figs.~\ref{fig:R_Ts_ma} and~\ref{fig:error_propagate} jointly examine Hour 12 as a representative mid-day hour. Fig.~\ref{fig:R_Ts_ma} shows the reserve-conditioned work-point family under different candidate reserve levels. As the candidate reserve increases, the optimizer lowers $T_s$ from $50.00^\circ\mathrm{C}$ to $35.00^\circ\mathrm{C}$ and increases $\dot m_a$ from $6.00~\mathrm{kg/s}$ to $17.63~\mathrm{kg/s}$. This work-point migration forms a scheduling-feasible reserve interval that extends to approximately $R_{\max}=55.5~\mathrm{kW}$, while the deployment-credible portion of this interval terminates earlier at $R_{\mathrm{cred}}\approx37.4~\mathrm{kW}$. The same representative hour also explains why the full scheduling-feasible interval cannot be regarded as credible. As shown in Fig.~\ref{fig:error_propagate}, at moderate reserve levels such as $R=15~\mathrm{kW}$ and $R=30~\mathrm{kW}$, the propagated uncertainty band remains within the risk-tightened comfort region. When the reserve level approaches $R_{\mathrm{cred}}$, the minimum thermal margin starts to become binding; if the reserve is further increased toward $R_{\max}$, the propagated uncertainty band crosses the risk-tightened comfort boundary. Therefore, the interval between $R_{\mathrm{cred}}$ and $R_{\max}$ remains feasible in the supervisory scheduling problem, but does not satisfy the deployment-oriented credibility condition.
	
	\begin{figure}[t]
		\centering
		\includegraphics[width=0.9\columnwidth]{fig10_R_Ts_ma.pdf}
		\caption{Illustration of the reserve-conditioned work-point family for a representative hour. The feasible domain extends to $R_{\max}$, while the credible domain terminates earlier at $R_{\mathrm{cred}}$.}
		\label{fig:R_Ts_ma}
	\end{figure}
	
	\begin{figure}[htbp]
		\centering
		\includegraphics[width=0.9\columnwidth]{fig11_error.pdf}
		\caption{Intra-hour thermal-risk propagation for a representative hour under different reserve levels. The dashed lines denote the risk-tightened comfort bounds, and the shaded region denotes the propagated uncertainty band.}
		\label{fig:error_propagate}
	\end{figure}
	
	These results demonstrate the transition from geometric headroom to scheduling feasibility and then to credible deliverability. Work-point reconstruction expands the reserve domain by relocating the water--air operating point, whereas deployment-oriented credibility screening contracts this domain according to intra-hour thermal risk. The gap between $R_{\max}(h)$ and $R_{\mathrm{cred}}(h)$ varies with the reconstructed work point, the regulation-induced temperature deviation, and the available risk-tightened comfort margin. The resulting credible interval $0\le R\le R_{\mathrm{cred}}(h)$ serves as the effective reserve domain for the hourly cost-curve analysis in the next subsection.
	
	\subsection{Hourly Cost Curves Over the Credible Reserve Domain}
	
	Based on the credible supply function defined in Section~\ref{subsec:cost_supply_function}, this subsection evaluates the hourly holding-cost curves obtained from the reserve-conditioned scheduling results. The purpose is to examine how the holding cost varies with reserve level across different hours and how credibility screening affects the economically relevant portion of the feasible cost curve.
	
	Fig.~\ref{fig:cost_wall} provides a global view of the hourly reserve cost surface over the day. Two features can be observed. First, the cost structure is strongly hour-dependent. Some hours admit a relatively wide low-cost reserve region, whereas others exhibit a steep increase even at moderate reserve levels. Second, the cost growth is generally nonlinear. In many hours, $C_h(R)$ remains close to zero over the lower part of the reserve range and then rises rapidly as the reserve level approaches the upper end of the admissible domain. This behavior is consistent with the proposed work-point reconstruction mechanism: larger reserve commitments require stronger migration of the water--air operating point and may induce additional energy expenditure and downstream thermal recovery cost.
	
	\begin{figure}[htbp]
		\centering
		\includegraphics[width=0.8\columnwidth]{fig12_cost_wall.pdf}
		\caption{Hourly reserve cost surface over the credible reserve domain. The surface represents the incremental holding cost $C_h(R)$ associated with reserve level $R$ at hour $h$.}
		\label{fig:cost_wall}
	\end{figure}
	
	To further illustrate the role of credibility screening, Fig.~\ref{fig:cost_curve_of_H10_H12} compares the cost curves at Hour 10 and Hour 12. For both representative hours, the cost increases gradually at low reserve levels and becomes markedly steeper as the reserve approaches the scheduling-feasible boundary $R_{\max}$. At Hour 10, the scheduling-feasible reserve can be extended to $R_{\max}=65.8~\mathrm{kW}$, whereas the credible domain ends at $R_{\mathrm{cred}}=55.2~\mathrm{kW}$. The corresponding cost at the credible boundary is $\$9.35$, while the cost at the scheduling-feasible upper bound reaches $\$41.55$. Thus, the last $10.6~\mathrm{kW}$ of scheduling-feasible but non-credible reserve is associated with an additional holding cost of $\$32.20$.
	
	A similar pattern is observed at Hour 12. The credible boundary is $R_{\mathrm{cred}}=37.4~\mathrm{kW}$, compared with the scheduling-feasible upper bound $R_{\max}=55.5~\mathrm{kW}$. The corresponding cost increases from $\$9.74$ at $R_{\mathrm{cred}}$ to $\$33.8$ at $R_{\max}$. Therefore, the credibility screen removes $18.1~\mathrm{kW}$ of reserve from the scheduling-feasible interval and avoids a high-cost tail of $\$24.06$. These two cases indicate that credibility screening does not merely reduce the reported reserve quantity; it also identifies the economically relevant portion of the feasible cost curve.
	
	\begin{figure}[htbp]
		\centering
		\includegraphics[width=1\columnwidth]{fig13_cost_curve_of_H10_H12.pdf}
		\caption{Representative hourly cost curves at Hours 10 and 12. The blue shaded region denotes the credible reserve domain, and the orange shaded region denotes the feasible-but-not-credible region.}
		\label{fig:cost_curve_of_H10_H12}
	\end{figure}
	
	Overall, these results show that hourly reserve provision should be characterized by the pair $(R,C_h(R))$ over the credible domain $0\le R\le R_{\mathrm{cred}}(h)$, rather than by the scheduling-feasible upper bound alone. Work-point reconstruction determines how far the reserve domain can be expanded under supervisory feasibility, while credibility screening identifies the part of that domain that remains deliverable under intra-hour thermal uncertainty. The resulting credible cost curve therefore reports both deliverable reserve quantity and the corresponding holding cost, providing a more suitable reserve-supply representation for subsequent downstream aggregation and allocation.
	
	% \subsection{Sensitivity to Confidence Level}
	
	% \begin{figure}[!ht]
		% 	\centering
		% 	\includegraphics[width=1\columnwidth]{fig14_credible_vs_max_reserve_blue.png}
		% 	\caption{Sensitivity of the hourly credible reserve profile to the confidence level. The purple curve denotes the scheduling-feasible reserve profile, and the shaded regions denote the corresponding credible reserve domains under different values of $\beta$.}
		% 	\label{fig:credible_vs_max_reserve_blue}
		% \end{figure}
	
	% Figure~\ref{fig:credible_vs_max_reserve_blue} compares the hourly credible reserve profiles under different confidence levels. The baseline denotes the scheduling-feasible reserve profile, while the shaded bands represent the corresponding credible reserve domains for $\beta=0.80$, $0.85$, $0.90$, and $0.95$. As expected, the credible reserve decreases monotonically as the confidence level increases. This reduction is mild during hours with ample thermal margin, but becomes much more pronounced during the midday period, where the reserve-supporting work point is closer to the tightened comfort boundary and the same deployment uncertainty leads to a larger loss of credible capacity.
	
	% These results indicate that the confidence level acts as a direct risk-tuning parameter for reserve qualification. A lower $\beta$ admits a wider credible reserve domain, whereas a higher $\beta$ yields a more conservative offer curve. Importantly, the impact of $\beta$ is not uniform over the day: the sensitivity is strongest exactly in those hours where thermal risk is already binding. This confirms that confidence-level selection should be viewed as a structured trade-off between reserve quantity and delivery reliability, rather than as a simple uniform derating of the feasible reserve profile.
	
	
	
	\section{Conclusion}
	\label{sec:conclusion}
	
	This paper proposed a work-point reconstruction framework for credible reserve characterization of water--air coupled HVAC systems. The key insight is that the fan-side reserve baseline is not fixed, but can be reshaped by the supply-water temperature through water--air coupling at the coil. Based on this mechanism, the proposed method reconstructs reserve-supporting operating points, screens the schedulable reserve domain under risk-calibrated thermal constraints, and maps the credible domain into an incremental holding-cost curve. The case study shows that water-side coordination can expand the scheduling-feasible reserve range, while credible screening further identifies the practically deliverable reserve region and removes high-cost points near the feasible boundary. These results indicate that HVAC reserve capability should be described by credible reserve--cost relations rather than deterministic geometric headroom alone. Future work will extend the proposed characterization to multi-building aggregation and market-oriented reserve bidding.
	
	\section*{Acknowledgments}
	The authors would also like to express their sincere appreciation to Shandong Zuoyao Technology Co., Ltd. and CET Shandong Electronic Co., Ltd. for their valuable assistance and support in the field implementation and data acquisition of the ASHP-DHS system used in this study.
	
	%{\appendix[Proof of the Zonklar Equations]
		%Use $\backslash${\tt{appendix}} if you have a single appendix:
		%Do not use $\backslash${\tt{section}} anymore after $\backslash${\tt{appendix}}, only $\backslash${\tt{section*}}.
		%If you have multiple appendixes use $\backslash${\tt{appendices}} then use $\backslash${\tt{section}} to start each appendix.
		%You must declare a $\backslash${\tt{section}} before using any $\backslash${\tt{subsection}} or using $\backslash${\tt{label}} ($\backslash${\tt{appendices}} by itself
		% starts a section numbered zero.)}
	
	
	\bibliographystyle{IEEEtran}
	\bibliography{tse_bibliography}
	
	%{\appendices
		%\section*{Proof of the First Zonklar Equation}
		%Appendix one text goes here.
		% You can choose not to have a title for an appendix if you want by leaving the argument blank
		%\section*{Proof of the Second Zonklar Equation}
		%Appendix two text goes here.}
	
	
\end{document}


